% HyperSafe
\begin{abstract}
Safety alignment in large language models is fragile under fine-tuning, as even benign adaptation can increase harmful compliance.
Existing defenses rely on per-checkpoint interventions (either retraining or modifying the model weights) or on model-agnostic classifiers that do not capture model-specific shifts, forcing a trade-off between deployment cost and recipe-specific protection.
We propose HyperSafe, a framework that generates a model-specific safety classifier at inference time.
Our approach builds on the observation that fine-tuning induces characteristic shifts in intermediate representations, which we use as an activation fingerprint of the fine-tuned model.
A hypernetwork maps these fingerprints to the parameters of a \textit{Safe Side Network} (\ssn{}) that classifies prompts as safe or harmful.
At deployment, a single forward pass produces an \ssn{} from a small set of calibration prompts, without gradient updates or model modification.
We evaluate HyperSafe on two model families (Qwen2-7B and LLaMA-3-8B) across three safety benchmarks (BeaverTails, AdvBench, HEx-PHI).
HyperSafe reduces harmful response rates from 19--31\% to under 1\% on every held-out checkpoint, while keeping downstream task accuracy within one percentage point of the fine-tuned baseline on average.
\end{abstract}
